% Solution: Gibbons (1992) Problem 2.17. Input by \solitem. Check: check_T3.py (item 27).
\begin{exsol}{Gibbons Problem 2.17}{136--137}{Advanced}{4}
\asked An infinitely lived firm faces a sequence of one-period workers. A worker chooses effort (output $y$, cost
$c < y$) or not; the firm then chooses a wage, with no enforceable contracts; each worker observes the whole past.
Describe strategies for the firm and the workers forming an SPE in which every worker works, and give a necessary and
sufficient condition.

\method A relational contract through trigger strategies (Gibbons \S2.3.D, pp.~107--112, efficiency wages; the
firm is the long-lived player). The worker's outside option is $0$ and a wage cannot be negative, so a worker
works only if the wage he expects, net of the effort cost, is at least $0$. The worst the workers can do to the firm
is never work, which gives it $0$.

\solution
\textbf{Strategies.} Fix a wage $w^* \in [c, y)$.
\begin{itemize}
  \item \emph{Worker $t$:} work iff in every earlier period in which output was produced the firm paid $w^*$ (and, in
  period~1, always); otherwise shirk.
  \item \emph{Firm:} after output pay $w^*$ as long as it has always done so; after any failure to pay $w^*$
  following output, pay $0$ for ever.
\end{itemize}
\textbf{Worker.} If he works he expects $w^*$ and gets $w^* - c \ge 0$, the value of his outside option, so working
is optimal. After the firm has once cheated, the worker expects wage $0$ and shirks (working would cost $c > 0$ for
nothing). That is a best reply.

\textbf{Firm.} On the path it gets $y - w^*$ per period: $\frac{y - w^*}{1-\delta}$. If it pays $0$ after output it keeps $y$
once and then no worker ever works again ($0$). It pays iff
\[
  \frac{y - w^*}{1-\delta} \ge y \iff w^* \le \delta y .
\]
After a deviation the firm pays $0$ and no output is produced, which is a stage best reply (a wage after zero output
costs it), so the punishment phase is itself an equilibrium.

\textbf{Condition.} A wage $w^*$ must satisfy $c \le w^* \le \delta y$, so the equilibrium exists iff
\[
  \delta \ge \frac{c}{y} .
\]
\textbf{Necessity.} In any SPE in which workers work, the equilibrium wage must be at least $c$ (workers can
refuse). The firm can always deviate to $0$, and the harshest punishment it can face is zero profit for ever,
because workers can always shirk. So the equilibrium wage is at most $\delta y$, which is impossible if $\delta <
\frac{c}{y}$. Hence $\delta \ge \frac cy$ is also necessary.

\interp The firm honours an unenforceable wage promise because reneging ends the relationship. The wage that
induces effort ($c$) is the minimum a worker accepts, and the bound $\delta y$ is the most the firm will pay
without being tempted to renege. The relationship is possible iff effort is cheap relative to the firm's patience
($\frac cy \le \delta$): output must exceed the effort cost with enough margin, discounted.

\verify \texttt{check\_T3.py} (item~27) checks, for $(y, c) \in \{(10,4), (10,9), (5,1)\}$ and all
$\delta = \frac k{100}$, that a wage in $[c, \delta y]$ exists iff $\delta \ge c/y$, that every such wage
satisfies the firm's constraint, and that when $\delta < c/y$ no wage $\ge c$ does. Gibbons's unofficial manual
is not copied.

\mistakes
\begin{itemize}
  \item Specifying the firm's strategy only on the path. SPE requires what happens after a failure to pay.
  \item Setting the wage at $y$ or at $y - c$: the cap is $\delta y$, and a wage above $y$ would not be paid.
\end{itemize}

\answer{Workers work iff the firm has always paid $w^*$ after output; the firm pays $w^*$ as long as it always has,
and $0$ for ever otherwise, with $c \le w^* \le \delta y$. The firm complies iff $\frac{y - w^*}{1-\delta} \ge y$, and the
worker iff $w^* \ge c$. Such $w^*$ exists, and the SPE exists, iff $\delta \ge \dfrac{c}{y}$ (also necessary).}
\end{exsol}
